% !TEX program = pdflatex
% 结构来源（BMC Bioinformatics Research article）：
% https://link.springer.com/journal/12859/submission-guidelines/research-article
% 官方模板下载：
% https://www.springernature.com/gp/authors/campaigns/latex-author-support
% 将官网 ZIP 中 sn-jnl.cls 和 sn-vancouver-num.bst 放入本目录后，
% 自动使用官方类；未放入时用 article 生成可编译的工作稿。
% article 回退版本不是官方投稿样式；投稿前必须使用官网 ZIP。
% 编译：pdflatex main; bibtex main; pdflatex main; pdflatex main
% 本稿只提供填表骨架；不自动确定任务、指标、第三模态或算法细节。

\newif\ifSNtemplate
\IfFileExists{sn-jnl.cls}{%
  \IfFileExists{sn-vancouver-num.bst}{\SNtemplatetrue}{\SNtemplatefalse}%
}{\SNtemplatefalse}
\ifSNtemplate
  \documentclass[pdflatex,sn-vancouver-num]{sn-jnl}
\else
  \documentclass[11pt,a4paper]{article}
  \usepackage[margin=25mm]{geometry}
  \usepackage[numbers,sort&compress]{natbib}
  \usepackage{hyperref}
  \bibliographystyle{unsrtnat}
\fi

\usepackage[T1]{fontenc}
\usepackage{amsmath}
\usepackage{xcolor}
\usepackage{graphicx}
\usepackage{booktabs}
% 新增 Table 1（tab:prototype-compensation）与方法公式所需：
% amssymb 提供 \mathbb；rotating 提供横向浮动体 sidewaystable（不缩放）。
\usepackage{amssymb}
\usepackage{rotating}
% ### Bug Post-Mortem
% - 现象: sn-jnl 标题页缺宏命令，缩放表格出现分组错误。
% - 根因: 遗漏 xcolor；resizebox 嵌套干扰类内 threeparttable。
% - 修复: 加载 xcolor，使用不缩放的 tabular*。
% - Prevention Rule: 官方类与 article 回退分支分别完整编译。
\hypersetup{hidelinks}
\setlength{\emergencystretch}{2em}

\newcommand{\TODO}[1]{\textbf{[TODO: #1]}}
% 每个 -- 均为待填单元格，不表示零或已确认的不适用。
% 若最终采用均值和标准差，可用 \result{均值}{标准差} 替换 --。
\newcommand{\result}[2]{\shortstack{#1\\$\pm$#2}}
\newcommand{\draftabstract}{%
  \textbf{Background:} \TODO{State the cancer prediction task, motivation,
  study objective, and model name.}\par
  \textbf{Results:} \TODO{Summarize the completed comparisons in
  Tables~\ref{tab:two} and~\ref{tab:three}, if Table~\ref{tab:three} is
  retained; insert only documented results.}\par
  \textbf{Conclusions:} \TODO{State conclusions supported by the completed
  comparisons. Do not insert a performance or novelty claim yet.}}
\newcommand{\draftkeywords}{Multimodal learning; cancer prediction;
  textual data; genomic data; multimodal fusion}

\ifSNtemplate
  \title[Multimodal cancer prediction]{TODO: Title of the multimodal
  cancer prediction study}
  \author*[1]{\fnm{TODO} \sur{Author}}
  \email{TODO}
  \affil*[1]{\orgname{TODO: institution},
    \orgaddress{TODO: full institutional address}}
  \abstract{\draftabstract}
  \keywords{\draftkeywords}
\else
  \title{TODO: Title of the multimodal cancer prediction study}
  \author{TODO: author names\\TODO: institutions and addresses\\
    TODO: corresponding author and email}
  \date{}
\fi

\begin{document}
\maketitle
\ifSNtemplate\else
  \begin{abstract}
    \draftabstract
  \end{abstract}
  \noindent\textbf{Keywords:} \draftkeywords
\fi

\section{Background}
\TODO{Introduce the clinical prediction task and the motivation for
combining textual and genomic information.}

The Cancer Genome Atlas (TCGA) provides the data-source context for this
draft~\cite{tcga2013pan_cancer}.
\TODO{Describe the specific data used and the research question.}

The baseline reporting framework uses Table 1 of Song et
al.~\cite{song2026cancer_type_aware} as the planned source of published
comparison values.
\TODO{Add the study rationale and a factual description of Ours once
the method has been specified.}

\section{Methods}
\label{sec:methods}

\subsection{Data and input checklist}
\begin{itemize}
  \item Cohort columns: BLCA, BRCA, GBMLGG, LUAD, UCEC, CRAD, HNSC,
  and STAD. \TODO{Record cohort definitions, sample counts, inclusion
  criteria, and data versions.}
  \item Two-modality inputs: text and genomics.
  \TODO{Record data fields, preprocessing, and feature representations.}
  \item Optional three-modality inputs: text, genomics, and
  \TODO{third modality and its representation}.
  \item Prediction target: \TODO{survival or classification, outcome
  definition, labels, and prediction time point if applicable}.
\end{itemize}

\subsection{Baseline comparison checklist}
\begin{itemize}
  \item Table~\ref{tab:two}: Cross-Attn Fusion, Early Fusion,
  Late Fusion, and Ours.
  \item Table~\ref{tab:three}, if retained: LIMOE, MAGGate, MulT,
  TF, Cross-Attn Fusion, Early Fusion, Late Fusion, and Ours.
  \item Ours: \TODO{model name, existing implementation identifier,
  and a short factual method description}.
\end{itemize}

Bibliographic slots are reserved for Cross-Attn Fusion,
Early Fusion, and Late Fusion~\cite{todo_cross_attention_fusion,
todo_early_fusion,todo_late_fusion}, and for LIMOE, MAGGate, and
Transformer Fusion (TF)~\cite{todo_limoe,todo_maggate,
todo_transformer_fusion}.
MulT has a candidate original reference~\cite{tsai2019mult}.
\TODO{Resolve each TODO-CITE and confirm the mapping from each
table label to its cited method.}

\subsection{Evaluation and filling checklist}
\begin{itemize}
  \item Validation for Ours: five-fold evaluation.
  \TODO{Record split files, fold identifiers, and their relationship
  to the source study.}
  \item Per-cohort reporting: fill the eight cohort columns in
  Table~\ref{tab:two} and, if retained, Table~\ref{tab:three}.
  \item Metric: \TODO{full metric name, abbreviation, implementation,
  and whether higher or lower values are better}.
  \item Summary within each cohort: \TODO{aggregation across folds
  and the uncertainty measure reported beside the point estimate}.
  \item Overall: \TODO{aggregation definition, weighting, and
  source or calculation record for each row}.
  \item Ours result records: \TODO{paths to existing per-fold results,
  configurations, and the checkpoint-selection rule}.
\end{itemize}

\subsection{Provenance of comparison values}
Non-Ours values in Tables~\ref{tab:two} and~\ref{tab:three} are planned
to be transcribed from Table 1 of Song et
al.~\cite{song2026cancer_type_aware}. These are literature-reported
results; they are not experiments reproduced by the authors of this
study. Ours rows are reserved for the authors' own experimental results.

\TODO{Record the source row, column, and setting for each transcribed
value. Record the calculation used for any derived Overall entry.}
All numerical entries remain unfilled in this draft.

% =====================================================================
% 新增小节（4.progress 交接）：原型引导补偿方法的公式（λ 统一表述）。
% 溯源（仅注释，不渲染）：协议
%   patient-fixed-padmask-v2-infer-rules1to5-v1-luad-m1-force-v1
% 公式是已执行计算链的数学化重述；每式上方的 "code:" 注释给出代码落点。
% 实现中的分场景路由在此写成 λ_{c,g} 与 γ_{c,g}，与实现逐格等价：
%   均值填补路由 ≡ γ=1、λ=0（pooling 权重 1）；不补偿路由 ≡ γ=0；
%   检索路由 ≡ γ=1、λ=1（pooling 权重 w）。
% 路径前缀：I01 = experiments/I01_patient_retrieval；
%   R2/R3/R4 = I01/knowledge/data-analysis-from-grok/{2.progress/result,
%   3.progress/result,4.progress}；SRC = src/trimodalsurv。
% 本小节范围为五癌、三模态、五 seed，与上文八癌/两模态/five-fold 占位清单无关。
% =====================================================================
\subsection{Prototype-guided compensation for missing modalities}
\label{sec:prototype-compensation}

\paragraph{Scope.}
This subsection specifies the inference-time method that is evaluated as the
third strategy in Table~\ref{tab:prototype-compensation}. It is a
non-parametric compensation guided by the cached whole-slide image (WSI)
prototypes of each patient: a missing modality is filled with features taken
from the most similar training patient, shrunk towards the training mean by a
setting-specific coefficient, and compensation can be switched off for a
setting. Nothing is trained in this study. The survival networks are frozen
(one previously trained NPJ-C model per cohort and seed, 25 in total), and the
method involves no learnable prototypes, no learned compensation module, no
moving-average update, no contrastive or alignment loss and no additional
training objective. The experiment covers five TCGA cohorts (BLCA, BRCA, LGG,
LUAD and UCEC), three modalities and five seeds; the seeds index independently
trained networks and are not cross-validation folds. The equations below
restate the executed computation and are not a derivation of optimality.

\paragraph{Modality availability.}
Let $c$ denote a cohort, $i$ a query patient, $j$ a candidate donor,
$m\in\{W,R,T\}$ the WSI, RNA and text modalities, $g$ an evaluation setting
and $\mathcal{I}_c$ the training patients of cohort $c$. The cached feature
array of modality $m$ before projection is $\mathbf{X}_i^{m}$, of size
$128\times1536$ for the WSI, $2048\times256$ for RNA and $200\times768$ for
text. With $n_i^{m}\in\{0,1\}$ indicating that modality $m$ is naturally
available and $b_g^{m}\in\{0,1\}$ the artificial mask of setting $g$, the
availability seen by the model is
% code: I01/evaluate.py: batch（natural & 非遮挡）；R2/inference.py: prepare_batch 的 original/natural/fusion 掩码
\begin{equation}
  v_{i,g}^{m}=n_i^{m}\,b_g^{m},
  \qquad
  \bigl(b_g^{W},b_g^{R},b_g^{T}\bigr)=
  \begin{cases}
    (1,1,1), & g=\texttt{none},\\
    (1,0,1), & g=\texttt{rna\_100},\\
    (1,1,0), & g=\texttt{text\_100},\\
    (1,0,0), & g=\texttt{both\_100}.
  \end{cases}
  \label{eq:pr-visible}
\end{equation}
By Eq.~\eqref{eq:pr-visible} the WSI is never masked, and the setting
\texttt{none} adds no artificial mask: naturally missing modalities remain
missing, so this setting is not a complete-data condition. The order
$(W,R,T)$ is notational; the implementation orders the modality tokens as
image, text, RNA.

\paragraph{Retrieval key from valid WSI prototypes.}
The WSI cache $\mathbf{P}_i=\mathbf{X}_i^{W}\in\mathbb{R}^{128\times1536}$
holds at most 128 rows per patient, namely patient-level cluster centres of
patch embeddings or, when fewer rows are available, the available rows
followed by zero padding. The number 128 is therefore an upper bound, and not
every row is a genuine prototype. The cache is used as stored; no clustering
is repeated. Valid rows are identified on the uncentred cache and averaged,
% code: I01/model.py: valid_row_mask；R2/02_wsi_meanpool_key/model.py: _signature
\begin{equation}
  a_{ik}=\mathbf{1}\bigl[\mathbf{P}_{ik}\neq\mathbf{0}\bigr],
  \qquad
  \bar{\mathbf{p}}_i
  =\frac{\sum_{k=1}^{128}a_{ik}\,\mathbf{P}_{ik}}{\sum_{k=1}^{128}a_{ik}}
  \in\mathbb{R}^{1536},
  \label{eq:pr-mask}
\end{equation}
where the only admissible layout is a non-zero prefix followed by a zero
suffix. The mean in Eq.~\eqref{eq:pr-mask} excludes padding from retrieval.
The key is then centred with the training mean of the same cohort,
% code: I01/model.py: FixedPatientBank.__init__（self.mean，患者等权）；R2/02_wsi_meanpool_key/model.py: _signature
\begin{equation}
  \boldsymbol{\mu}_c^{W}
  =\frac{1}{|\mathcal{I}_c|}\sum_{j\in\mathcal{I}_c}\bar{\mathbf{p}}_j,
  \qquad
  \mathbf{q}_i=\bar{\mathbf{p}}_i-\boldsymbol{\mu}_c^{W}.
  \label{eq:pr-key}
\end{equation}
In Eq.~\eqref{eq:pr-key} every training patient has the same weight regardless
of its number of patches, and validation and test patients never contribute to
$\boldsymbol{\mu}_c^{W}$. The retrieval key $\mathbf{q}_i$ is a single
1536-dimensional vector; the row-aligned cosine over the flattened
$128\times1536$ array used by an earlier protocol is not part of this one.

\paragraph{Similarity, candidates and donor selection.}
The similarity between query $i$ and training patient $j$ is
% code: R2/02_wsi_meanpool_key/model.py: meanpool_wsi_score；R2/04_wsi_rna_joint_sim/model.py: token_mean_cosine, joint_similarity
\begin{equation}
  s_{ij}^{W}
  =\frac{\mathbf{q}_i^{\top}\mathbf{q}_j}
        {\lVert\mathbf{q}_i\rVert_2\,\lVert\mathbf{q}_j\rVert_2},
  \qquad
  s_{ij}^{R}
  =\frac{\bar{\mathbf{r}}_i^{\top}\bar{\mathbf{r}}_j}
        {\lVert\bar{\mathbf{r}}_i\rVert_2\,\lVert\bar{\mathbf{r}}_j\rVert_2},
  \qquad
  s_{ij}=s_{ij}^{W}+J_{i,g}\,\alpha\,s_{ij}^{R},
  \label{eq:pr-sim}
\end{equation}
where $\bar{\mathbf{r}}_i$ is the token mean of $\mathbf{X}_i^{R}$ without
centring and $\alpha=1$. The indicator $J_{i,g}$ equals one only if
$g=\texttt{text\_100}$, features are retrieved in that setting
($\gamma_{c,g}=1$ and $\lambda_{c,g}=1$, both defined below) and the patient's
own RNA is naturally available; otherwise $J_{i,g}=0$ and no RNA is read. In
Eq.~\eqref{eq:pr-sim} a centred key norm or an RNA norm not exceeding
$10^{-12}$ raises an error instead of being regularised.
With $\mathcal{M}_{i,g}=\{m\in\{R,T\}:v_{i,g}^{m}=0\}$ the modalities to be
compensated, the admissible donors are
% code: R2/inference.py: _custom_retrieve（候选过滤；加速路径为等价查询，不是新方法）
\begin{equation}
  \mathcal{D}_{i,g}
  =\Bigl\{j\in\mathcal{I}_c:\ j\neq i,\ \
  n_j^{m}=1\ \ \forall m\in\mathcal{M}_{i,g},\ \
  \bigl(J_{i,g}=1\Rightarrow n_j^{R}=n_j^{T}=1\bigr)\Bigr\}.
  \label{eq:pr-cand}
\end{equation}
Equation~\eqref{eq:pr-cand} restricts donors to training patients of the same
cohort that naturally possess every modality the query lacks, and it excludes
the query itself. When both RNA and text are missing, a donor must possess
both, so that one donor supplies both modalities. Validation and test patients
never enter the patient bank, survival times and censoring indicators are not
used for retrieval, and an empty candidate set raises an error. The donor is
the most similar admissible patient,
% code: R2/inference.py: _custom_retrieve（patient ID 已排序，argmax 取最前）
\begin{equation}
  j^{\star}
  =\operatorname*{arg\,max}_{j\in\mathcal{D}_{i,g}}s_{ij},
  \label{eq:pr-top1}
\end{equation}
and ties in Eq.~\eqref{eq:pr-top1} are resolved by taking the first patient in
sorted identifier order, which makes retrieval deterministic.

\paragraph{Compensated value.}
Let $\mathcal{I}_c^{m}=\{j\in\mathcal{I}_c:n_j^{m}=1\}$. The training mean of
modality $m$ and the compensated value are
% code: I01/model.py: FixedPatientBank.__init__（self.means，逐元素均值）；R2/03_shrinkage_lambda/model.py: shrink_feature
% code（λ_{c,g}、γ_{c,g} 取值的实现落点）: R4/luad_m1_force/model.py: route_patient；R2/01_text100_skip_compensate/model.py: should_skip_compensation；R3/strategy0_inference.py 的批内路由
\begin{equation}
  \boldsymbol{\mu}_c^{m}
  =\frac{1}{|\mathcal{I}_c^{m}|}\sum_{j\in\mathcal{I}_c^{m}}\mathbf{X}_j^{m},
  \qquad
  \widetilde{\mathbf{X}}_i^{m}
  =\boldsymbol{\mu}_c^{m}
   +\lambda_{c,g}\bigl(\mathbf{X}_{j^{\star}}^{m}-\boldsymbol{\mu}_c^{m}\bigr),
  \qquad \lambda_{c,g}\in\{0,1\}.
  \label{eq:pr-shrink}
\end{equation}
The mean in Eq.~\eqref{eq:pr-shrink} is element-wise over complete feature
arrays of identical shape. The WSI prototypes act only as the retrieval key:
the value returned is the donor's complete pre-projection feature array of the
target modality, and cluster centres are never substituted for RNA or text.
The coefficient $\lambda_{c,g}$ is specific to cohort $c$ and setting $g$.
With $\lambda_{c,g}=1$ the compensated value equals the donor's original
features, selected by Eqs.~\eqref{eq:pr-sim}--\eqref{eq:pr-top1}; with
$\lambda_{c,g}=0$ the donor term vanishes, the value reduces to the training
mean and no retrieval is performed. In addition, a binary gate
$\gamma_{c,g}\in\{0,1\}$ determines whether compensation is applied at all in
cohort $c$ and setting $g$: with $\gamma_{c,g}=0$ the missing modalities
remain masked. Both reference strategies are special cases of this
formulation, namely no compensation ($\gamma_{c,g}=0$ throughout) and mean
imputation ($\gamma_{c,g}=1$ and $\lambda_{c,g}=0$ throughout). The constants
$\alpha$ and $w$ (defined below) and the default $\lambda_{c,g}=1$ were
selected on the validation split.
\TODO{Author to add: values of $\lambda_{c,g}$ and $\gamma_{c,g}$ for each
cohort and setting (needed to reproduce the identical entries of
Table~\ref{tab:prototype-compensation}).}

\paragraph{Filling and fusion mask.}
The model input and the fusion validity flag are
% code: R2/inference.py: prepare_batch；I01/model.py: FixedPatientBank.compensate
\begin{equation}
  \mathbf{X}_{i,\mathrm{in}}^{m}=
  \begin{cases}
    \mathbf{X}_i^{m}, & v_{i,g}^{m}=1\ \text{or}\ \gamma_{c,g}=0,\\
    \widetilde{\mathbf{X}}_i^{m}, & v_{i,g}^{m}=0\ \text{and}\ \gamma_{c,g}=1,
  \end{cases}
  \label{eq:pr-fill}
\end{equation}
% code: R2/inference.py: prepare_batch 的 fusion/imputed 掩码；I01/model.py: FixedPatientBank.compensate 返回的 masks
\begin{equation}
  f_i^{m}
  =\max\Bigl\{v_{i,g}^{m},\
   \gamma_{c,g}\,\mathbf{1}\bigl[m\in\mathcal{M}_{i,g}\bigr]\Bigr\}.
  \label{eq:pr-flag}
\end{equation}
By Eq.~\eqref{eq:pr-fill} features that the patient already has are never
replaced, and when both modalities are filled with $\lambda_{c,g}=1$ they come
from the same donor $j^{\star}$. With $\gamma_{c,g}=0$ the missing slot keeps
a placeholder that Eq.~\eqref{eq:pr-flag} marks as invalid, so it is excluded
from fusion; with $\gamma_{c,g}=1$ the filled slot is marked valid. Means,
similarities and the compensated value are computed in double precision and
cast back to the input precision of the network.

\paragraph{Frozen encoder and fusion.}
Each modality enters the frozen network through its original interface,
% code: SRC/models/npjc.py: NPJC.encode_patient_modalities（mean(dim=1) → m_projector）与 modality_embed
\begin{equation}
  \mathbf{z}_i^{m}
  =\phi_m\bigl(\operatorname{mean}_{\mathrm{token}}
   (\mathbf{X}_{i,\mathrm{in}}^{m})\bigr),
  \qquad
  \mathbf{e}_i^{m}=\mathbf{z}_i^{m}+\mathbf{d}_m,
  \qquad
  \phi_m=\operatorname{Dropout}\circ\operatorname{ReLU}\circ
         \operatorname{Linear}_m,
  \label{eq:pr-proj}
\end{equation}
where the projector $\phi_m$ and the modality embedding $\mathbf{d}_m$ are
frozen parameters and dropout is inactive in evaluation mode.
Equation~\eqref{eq:pr-proj} implies that a retrieved array is first averaged
over its tokens and then projected, yielding one vector per modality; donor
tokens are not passed to the fusion Transformer individually. The patient's
own WSI is encoded as the mean over all 128 cached rows, including zero
padding, which is the historical behaviour of the frozen network; the
valid-row mean of Eq.~\eqref{eq:pr-mask} is used for retrieval only. The
three modality vectors are fused by the frozen Transformer encoder,
% code: R2/05_imputed_token_downweight/model.py: weighted_forward（backbone, src_key_padding_mask=~valid）
\begin{equation}
  \bigl(\mathbf{h}_i^{W},\mathbf{h}_i^{R},\mathbf{h}_i^{T}\bigr)
  =\operatorname{Transformer}_{\theta}\Bigl(
   \bigl(\mathbf{e}_i^{W},\mathbf{e}_i^{R},\mathbf{e}_i^{T}\bigr);\
   \text{key-padding mask}=\neg\,\mathbf{f}_i\Bigr),
  \label{eq:pr-trans}
\end{equation}
with frozen parameters $\theta$. In Eq.~\eqref{eq:pr-trans} filled modalities
take part in attention as valid tokens. The output tokens are pooled with
fixed weights,
% code: R2/05_imputed_token_downweight/model.py: build_pool_weights, weighted_forward（clamp_min(1)）
\begin{equation}
  \omega_i^{m}=
  \begin{cases}
    0, & f_i^{m}=0,\\
    w, & v_{i,g}^{m}=0,\ f_i^{m}=1,\ \lambda_{c,g}=1,\\
    1, & \text{otherwise},
  \end{cases}
  \qquad
  \mathbf{u}_i
  =\frac{\sum_{m}\omega_i^{m}\,\mathbf{h}_i^{m}}
        {\max\bigl(1,\ \sum_{m}\omega_i^{m}\bigr)},
  \qquad w=0.5.
  \label{eq:pr-pool}
\end{equation}
In Eq.~\eqref{eq:pr-pool} the weight $w=0.5$ applies only to modalities filled
with retrieved features ($\lambda_{c,g}=1$). A modality filled with the
training mean ($\lambda_{c,g}=0$) has weight one, exactly as in the
mean-imputation reference, and a masked modality has weight zero. Because the
WSI is always
valid the denominator is at least one, and the maximum mirrors the clamp of
the implementation. The weight $w$ is a hand-set constant and not a learned
reliability gate, and it does not remove the influence that a filled token
exerts on the other tokens through attention in Eq.~\eqref{eq:pr-trans}.

\paragraph{Risk score and aggregation.}
The existing cohort-specific survival head $H_c$ yields four discrete-time
logits, from which the risk score is computed as
% code: SRC/evaluation/common.py: dual_risk_from_logits（第 2 个返回值，B 口径）；I01/evaluate.py: evaluate_unit
\begin{equation}
  \boldsymbol{\ell}_i=H_c(\mathbf{u}_i),
  \qquad
  \eta_{ik}=\operatorname{clip}\bigl(\sigma(\ell_{ik}),\,10^{-6},\,1-10^{-6}\bigr),
  \qquad
  S_{ik}=\prod_{\tau=1}^{k}\bigl(1-\eta_{i\tau}\bigr),
  \qquad
  R_i=-\sum_{k=1}^{4}S_{ik}.
  \label{eq:pr-risk}
\end{equation}
The concordance index (C-index) is computed from the risk $R_i$ of
Eq.~\eqref{eq:pr-risk}, the survival time and the event indicator with the
same implementation for all three strategies; no time bins, loss weights or
training hyperparameters are changed. With $C_{c,g,a,s}$ the test C-index of
cohort $c$, setting $g$, strategy $a$ and seed $s$, the quantities reported in
Table~\ref{tab:prototype-compensation} are
% code: R3/strategy0_reporting.py: summarize（R4/luad_m1_force/reporting.py 调用）
\begin{equation}
  \bar{C}_{c,g,a}=\frac{1}{5}\sum_{s\in\mathcal{S}}C_{c,g,a,s},
  \qquad
  O_{c,a}=\frac{1}{4}\sum_{g\in\mathcal{G}}\bar{C}_{c,g,a},
  \qquad
  \bar{C}_{a}=\frac{1}{5}\sum_{c}O_{c,a},
  \label{eq:pr-agg}
\end{equation}
where $\mathcal{S}=\{123,132,213,231,321\}$ are the seeds, $\mathcal{G}$
contains the four settings and $a$ ranges over no compensation (m0real), mean
imputation (m1) and the proposed prototype-guided compensation (strategy0). By
Eq.~\eqref{eq:pr-agg} every entry for a single setting averages five seeds,
the overall value $O_{c,a}$ of a cohort averages its 20 cells, and the
macro-average $\bar{C}_{a}$ averages all 100 cells with equal weights;
patients are not pooled across seeds and no seed is selected. The number of
cohorts in which the proposed method exceeds both references is
% code: R3/strategy0_reporting.py: summarize（cancer_wins，严格大于）
\begin{equation}
  W=\sum_{c}\mathbf{1}\bigl[\,O_{c,a^{\star}}>O_{c,\mathrm{m1}}\ \land\
     O_{c,a^{\star}}>O_{c,\mathrm{m0real}}\bigr],
  \label{eq:pr-wins}
\end{equation}
where $a^{\star}$ denotes the proposed method (strategy0). The comparisons in
Eq.~\eqref{eq:pr-wins} use unrounded values.

\section{Results}
\label{sec:results}

% =====================================================================
% 新增小节（4.progress 交接）：Table 1 与结果说明。
% 第三列 Proto. = JSON 中的 strategy0（协议 ID 见方法小节头注释，仅供仓库内溯源）。
% 数值唯一来源：4.progress/runs/test/complete.json 的 rows/c_index_b
% （300 条 = 5 癌 × 5 seed × 4 场景 × 3 臂；SHA-256 8604aae4…2260eb6）。
% 单元格 = 五 seed 未舍入均值；Overall = 同癌同方法 20 格均值；展示六位。
% 加粗按未舍入值比较，并列全部加粗。表格布局依据
% manuscript/消融实验/消融实验-prototype.excalidraw.svg。
% 本表须位于 tab:two 之前，才能自动编号为 Table 1。
% =====================================================================
\subsection{Compensation strategies under missing modalities}
\label{sec:prototype-compensation-results}

Table~\ref{tab:prototype-compensation} compares three inference-time
strategies for missing modalities on identical frozen networks: no
compensation, mean imputation, and the proposed prototype-guided compensation
defined in Section~\ref{sec:prototype-compensation}. Averaged with equal
weights over the five cohorts, five seeds and four settings
(Eq.~\eqref{eq:pr-agg}), the proposed method reached a macro-average C-index
of 0.640506, compared with 0.631824 for mean imputation and 0.630751 without
compensation. In the present evaluation the proposed method thus increased
the average discriminative performance. Values below are given in the order
proposed method, mean imputation, no compensation. The cohort-level overall
value was higher than both references in four of the five cohorts (BRCA, LGG,
LUAD and UCEC; $W=4$ in Eq.~\eqref{eq:pr-wins}), although the margin in LUAD
was small (0.587039 versus 0.587004 and 0.586206), and it was not higher in
BLCA (0.578982 versus 0.582191 and 0.583388). This count refers to the overall
averages and does not mean that the proposed method was best in every setting.
The larger differences occurred when both auxiliary modalities were masked, in
BRCA (0.640263 versus 0.614738 and 0.596026), UCEC (0.609348 versus 0.575072
and 0.576319) and LGG (0.665286 versus 0.635266 and 0.576190), and for UCEC in
the text-masked setting (0.609748 versus 0.577743 and 0.589272). When only RNA
was masked, or without additional masking, the three strategies were close,
and mean imputation or no compensation was highest in several cohorts. No
statistical test was performed.
\TODO{Author to add: limitation concerning how the LUAD setting was selected.}

\begin{sidewaystable}
\centering
\caption{Comparison of three compensation strategies across five cancers
under four modality-availability settings. Entries are test C-indices
averaged over five seeds; \emph{Overall} is the equally weighted mean of the
four settings of a cohort. Proto.\ denotes the proposed prototype-guided
compensation (Section~\ref{sec:prototype-compensation}).}
\label{tab:prototype-compensation}
\scriptsize
\setlength{\tabcolsep}{2pt}
\renewcommand{\arraystretch}{1.35}
\begin{tabular*}{\linewidth}{@{\extracolsep{\fill}}l*{15}{c}@{}}
\toprule
 & \multicolumn{3}{c}{BRCA} & \multicolumn{3}{c}{UCEC}
 & \multicolumn{3}{c}{LUAD} & \multicolumn{3}{c}{BLCA}
 & \multicolumn{3}{c}{LGG} \\
\cmidrule(lr){2-4}\cmidrule(lr){5-7}\cmidrule(lr){8-10}
\cmidrule(lr){11-13}\cmidrule(lr){14-16}
Scenario
 & No comp. & Mean & Proto.
 & No comp. & Mean & Proto.
 & No comp. & Mean & Proto.
 & No comp. & Mean & Proto.
 & No comp. & Mean & Proto. \\
\midrule
RNA masked & 0.683842 & 0.684751 & \textbf{0.685705} & 0.647229 & \textbf{0.655464} & 0.651413 & \textbf{0.580660} & 0.580239 & 0.580239 & \textbf{0.593546} & 0.589105 & 0.588168 & 0.767219 & \textbf{0.790959} & 0.786611 \\
Text masked & \textbf{0.655002} & 0.619121 & \textbf{0.655002} & 0.589272 & 0.577743 & \textbf{0.609748} & \textbf{0.596630} & 0.595156 & \textbf{0.596630} & \textbf{0.575989} & 0.574913 & \textbf{0.575989} & \textbf{0.705866} & 0.664113 & \textbf{0.705866} \\
Both masked & 0.596026 & 0.614738 & \textbf{0.640263} & 0.576319 & 0.575072 & \textbf{0.609348} & 0.585469 & \textbf{0.588803} & \textbf{0.588803} & 0.574046 & \textbf{0.574774} & 0.561797 & 0.576190 & 0.635266 & \textbf{0.665286} \\
None & \textbf{0.691359} & 0.686068 & 0.688112 & 0.651502 & \textbf{0.656888} & 0.650968 & 0.582064 & \textbf{0.583819} & 0.582485 & \textbf{0.589972} & \textbf{0.589972} & \textbf{0.589972} & 0.796825 & \textbf{0.799517} & 0.797723 \\
\midrule
Overall & 0.656557 & 0.651170 & \textbf{0.667270} & 0.616081 & 0.616292 & \textbf{0.630369} & 0.586206 & 0.587004 & \textbf{0.587039} & \textbf{0.583388} & 0.582191 & 0.578982 & 0.711525 & 0.722464 & \textbf{0.738872} \\
\bottomrule
\end{tabular*}
\par\smallskip
\begin{minipage}{\linewidth}
\scriptsize
Scenarios: RNA masked, Text masked and Both masked remove the RNA profile,
the pathology report or both for all test patients (\texttt{rna\_100},
\texttt{text\_100}, \texttt{both\_100}); None (\texttt{none}) adds no
artificial mask but retains naturally occurring missingness and is therefore
not a complete-data condition. Overall is an average and not a fifth
experimental scenario. Strategies: No comp., no compensation (m0real), missing
modalities are masked; Mean, mean imputation (m1), a missing modality is
replaced by the element-wise training mean of the same cohort; Proto., the
proposed prototype-guided compensation (strategy0;
Section~\ref{sec:prototype-compensation}). Bold marks the
highest unrounded value among the three strategies within a cohort and row;
ties are all bold. The three strategies use the same 25 frozen networks; no
model was retrained and no statistical test was performed.
\end{minipage}
\end{sidewaystable}

\subsection{Two-modality comparison}
\TODO{Fill Table~\ref{tab:two} from the designated sources and the existing
Ours results, then describe the observed per-cohort and Overall comparisons.}

\begin{table}[htbp]
\centering
\caption{Two-modality comparison using text and genomics.
Metric and uncertainty summary: TODO.}
\label{tab:two}
\footnotesize
\setlength{\tabcolsep}{3pt}
\renewcommand{\arraystretch}{1.3}
\begin{tabular*}{\textwidth}{@{\extracolsep{\fill}}l*{9}{c}@{}}
\toprule
Method & BLCA & BRCA & GBMLGG & LUAD & UCEC & CRAD & HNSC & STAD & Overall \\
\midrule
Cross-Attn Fusion & -- & -- & -- & -- & -- & -- & -- & -- & -- \\
Early Fusion     & -- & -- & -- & -- & -- & -- & -- & -- & -- \\
Late Fusion      & -- & -- & -- & -- & -- & -- & -- & -- & -- \\
\textbf{Ours}    & -- & -- & -- & -- & -- & -- & -- & -- & -- \\
\bottomrule
\end{tabular*}
\par\smallskip
\begin{minipage}{\textwidth}
\footnotesize
Non-Ours values are planned to be taken from Table 1 of
Song et al.~\cite{song2026cancer_type_aware}, not reproduced in this
study. Ours values will be supplied from the authors' own experiments.
The symbol -- denotes an unfilled entry. Overall: TODO, including
whether it is transcribed or calculated.
\end{minipage}
\end{table}

\subsection{Three-modality comparison: optional}
\TODO{Confirm whether Table~\ref{tab:three} is retained and specify the
third modality. If retained, fill its entries and describe the comparisons.}

\begin{table}[htbp]
\centering
\caption{Optional three-modality comparison using text, genomics,
and a third modality (TODO). Metric and uncertainty summary: TODO.}
\label{tab:three}
\footnotesize
\setlength{\tabcolsep}{3pt}
\renewcommand{\arraystretch}{1.3}
\begin{tabular*}{\textwidth}{@{\extracolsep{\fill}}l*{9}{c}@{}}
\toprule
Method & BLCA & BRCA & GBMLGG & LUAD & UCEC & CRAD & HNSC & STAD & Overall \\
\midrule
LIMOE            & -- & -- & -- & -- & -- & -- & -- & -- & -- \\
MAGGate          & -- & -- & -- & -- & -- & -- & -- & -- & -- \\
MulT             & -- & -- & -- & -- & -- & -- & -- & -- & -- \\
TF               & -- & -- & -- & -- & -- & -- & -- & -- & -- \\
Cross-Attn Fusion & -- & -- & -- & -- & -- & -- & -- & -- & -- \\
Early Fusion     & -- & -- & -- & -- & -- & -- & -- & -- & -- \\
Late Fusion      & -- & -- & -- & -- & -- & -- & -- & -- & -- \\
\textbf{Ours}    & -- & -- & -- & -- & -- & -- & -- & -- & -- \\
\bottomrule
\end{tabular*}
\par\smallskip
\begin{minipage}{\textwidth}
\footnotesize
Non-Ours values are planned to be taken from Table 1 of
Song et al.~\cite{song2026cancer_type_aware}, not reproduced in this
study. Ours values will be supplied from the authors' own experiments.
TF denotes Transformer Fusion. The symbol -- denotes an unfilled entry.
Overall: TODO, including whether it is transcribed or calculated.
\end{minipage}
\end{table}

\section{Discussion}
\TODO{Interpret the completed comparisons without inventing a
mechanistic explanation or claiming an unmeasured improvement.}

\TODO{Discuss the findings in relation to the cited methods,
distinguishing literature-reported values from the authors' results.}

\TODO{State the limitations supported by the final study design
and completed evidence.}

\section{Conclusions}
\TODO{Summarize the findings and their scope after the result tables
have been completed.}

\section*{Abbreviations}
\begin{description}
  \item[TCGA] The Cancer Genome Atlas.
  \item[TF] Transformer Fusion.
  \item[Cross-Attn] Cross-Attention.
\end{description}
\TODO{Add the full names of the reported metric, cohort labels,
and remaining method abbreviations; confirm the definition of CRAD.}

\section*{Declarations}
\subsection*{Ethics approval and consent to participate}
\TODO{Author-confirmed ethics and consent statement, including
approval, exemption, or applicability details.}

\subsection*{Consent for publication}
\TODO{Author-confirmed consent-for-publication statement.}

\subsection*{Availability of data and materials}
\TODO{Dataset access details and links to available code, split
files, and result records. State any actual access restrictions.}

\subsection*{Competing interests}
\TODO{Author-confirmed declaration of competing interests.}

\subsection*{Funding}
\TODO{Funding bodies, grant identifiers, and funder roles,
or an author-confirmed statement of no funding.}

\subsection*{Authors' contributions}
\TODO{Author names or initials and their actual contributions.}

\subsection*{Acknowledgements}
\TODO{Acknowledgements and any required disclosure of writing
assistance, after author confirmation.}

\clearpage
\bibliography{references}
\end{document}
